\subsection{Solution to Problem 33: Two Masses with Pulley}

\begin{solution}
The larger mass ($4m$) moves downward with speed $v$, so the smaller mass ($2m$) moves upward with the same speed. Each mass experiences air resistance of magnitude $kv$ opposing its motion.

\textbf{(i)} Take downward as positive for the $4m$ mass. The forces are weight $4mg$ (down), tension $T$ (up), and resistance $kv$ (up), so Newton's second law gives
\[
4mg - T - kv = 4m\,\frac{dv}{dt}.
\]

For the $2m$ mass, take upward as positive. The forces are tension $T$ (up), weight $2mg$ (down), and resistance $kv$ (down), so
\[
T - 2mg - kv = 2m\,\frac{dv}{dt}.
\]

The string is light and the pulley is smooth, so the tension is the same on both masses. Adding the two equations eliminates $T$:
\[
(4mg - T - kv) + (T - 2mg - kv) = 6m\,\frac{dv}{dt},
\]
\[
2mg - 2kv = 6m\,\frac{dv}{dt}.
\]

Dividing by $6m$,
\[
\frac{dv}{dt} = \frac{2mg - 2kv}{6m} = \frac{gm - kv}{3m}.
\]

\textbf{(ii)} The motion satisfies
\[
\frac{dv}{dt} = \frac{gm - kv}{3m}.
\]
Separate variables:
\[
\frac{dv}{gm - kv} = \frac{dt}{3m}.
\]
Integrate both sides:
\[
-\frac{1}{k}\ln|gm - kv| = \frac{t}{3m} + C.
\]
The masses start from rest, so $v = 0$ when $t = 0$. Substituting gives
\[
-\frac{1}{k}\ln(gm) = C.
\]
Hence
\[
-\frac{1}{k}\ln|gm - kv| = \frac{t}{3m} - \frac{1}{k}\ln(gm).
\]
Rearranging,
\[
\ln\left|\frac{gm - kv}{gm}\right| = -\frac{kt}{3m}.
\]
We are given $v < \dfrac{gm}{k}$, so $gm - kv > 0$ and we may drop the absolute value:
\[
\frac{gm - kv}{gm} = e^{-kt/(3m)}.
\]
Solving for $v$,
\[
v = \frac{gm}{k}\left(1 - e^{-kt/(3m)}\right).
\]
When $t = \dfrac{3m}{k}\ln 2$,
\[
v = \frac{gm}{k}\left(1 - e^{-\ln 2}\right)
= \frac{gm}{k}\left(1 - \frac{1}{2}\right)
= \frac{gm}{2k}.
\]
\end{solution}

\begin{takeaways}
\begin{itemize}
    \item Newton's second law: Apply to each mass separately with directions chosen to match the motion
    \item System of equations: Add the equations to eliminate tension
    \item Separable ODE: Integrate and use the initial condition $v(0) = 0$ to find velocity as a function of time
\end{itemize}
\end{takeaways}
